\exercises

\tierunderstand

\begin{bt}
  \item Dẫn xuất ở \eqref{eq:ch07-phuong-sai} dùng giả thiết độc lập đúng hai
    lần, ở hai chỗ khác nhau. Chỉ ra cả hai. Rồi giả sử \(q_i\) và \(k_i\)
    tương quan với nhau, hệ số \(\rho\) như nhau ở mọi chiều: viết lại kỳ vọng
    và phương sai của \(\vect{q} \cdot \vect{k}\), và nói \(\sqrt{\dk}\) còn là
    hệ số chuẩn hóa đúng nữa không.

  \item Mục~\ref{sec:ch07-chia-can} đo được rằng ở lúc khởi tạo, phép chia cho
    \(\sqrt{\dk}\) chia quá tay khoảng ba lần, và kết luận rằng chuyện đó không
    làm hỏng lập luận. Viết ra lý do bằng lời của bạn. Rồi làm phần ngược lại,
    khó hơn: dựng một tình huống trong đó một hệ số hằng \emph{sai ba lần} thật
    sự làm hỏng mô hình, và nói tình huống ấy khác tình huống ở đây chỗ nào.

  \item Chứng minh trực tiếp từ \eqref{eq:ch07-attention} rằng attention giao
    hoán với phép hoán vị: nếu \(\mat{P}\) là ma trận hoán vị thì
    \(\attention(\mat{P}\mat{Q}, \mat{P}\mat{K}, \mat{P}\mat{V}) = \mat{P}\,
    \attention(\mat{Q}, \mat{K}, \mat{V})\). Chỉ rõ bước nào trong chứng minh
    dùng tới việc softmax lấy trên chiều key. Rồi nói vì sao cộng
    \eqref{eq:ch07-pe} vào phá được đẳng thức ấy.

  \item \eqref{eq:ch07-tich-vo-huong-pe} cho dạng đóng
    \(\sum_i \cos(k\omega_i)\), và mục~\ref{sec:ch07-vi-tri} đo được chỗ tổng ấy
    quay đầu: \(k = 5\) ở \(\dmodel = 32\), \(k = 6\) ở 64, \(k = 44\) ở 512.
    Từ dạng đóng, giải thích vì sao chỗ quay đầu dịch sang phải khi \(\dmodel\)
    lớn lên. Rồi dự đoán nó dịch đi đâu nếu giữ \(\dmodel = 64\) mà đổi hằng số
    10000 thành 100.

  \item Số tham số của một tầng multi-head không đổi theo \(h\). Viết ra lý do
    bằng một câu. Rồi nói ba thứ \emph{có} đổi theo \(h\), và với mỗi thứ nói
    nó đổi theo chiều nào.

  \item Mục~\ref{sec:ch07-che} đo được mô hình bỏ mask có mất mát huấn luyện
    0.0434, thấp hơn 0.0639 của mô hình đúng, và khớp đúng cả câu 0.0000. Giải
    thích cả hai con số. Rồi trả lời câu đắt hơn: nếu bạn chỉ được nhìn các đại
    lượng tính trong lúc huấn luyện, không được giải mã câu nào, bạn theo dõi
    cái gì để phát hiện lỗi này?

  \item Ở 3 epoch, mô hình bỏ mask có mất mát 0.6731 còn mô hình có mask được
    0.5443, tức ngược dấu so với ở 14 epoch. Giải thích chỗ đổi dấu. Rồi nói
    quan sát ấy nói gì về việc chọn thời điểm dừng cho một phép đo so sánh
    hai cấu hình.

  \item \eqref{eq:ch07-post-ln} đặt layer norm sau phép cộng. Viết đường đi của
    gradient qua một chồng \(N\) tầng con dưới thứ tự ấy, rồi viết lại dưới thứ
    tự ngược, \(\vect{x} + \mathrm{Sublayer}(\layernorm(\vect{x}))\). Chỉ ra
    chỗ khác nhau, và nói nó liên quan thế nào tới đường đi ở
    \eqref{eq:ch06-duong-moi} và tới vòng quay lỗi của
    chương~\ref{ch:lstm}.

  \item Mục~\ref{sec:ch07-corpus} đọc kiểu lỗi còn lại là \enquote{đúng từ, sai
    cụm} và quy nó cho thiếu thiên kiến quy nạp. Nêu một cách giải thích khác
    cho cùng kiểu lỗi ấy, rồi nói một phép đo tách được hai cách giải thích.
\end{bt}

\tierapply

Trạng thái xuất phát cho cả năm bài: clone repo companion, đứng ở thư mục gốc
của nó, và checkout đúng tag của chương này.

\begin{minted}{console}
$ git clone https://github.com/Giang-Dang/rnn-to-transformer-lab.git
$ cd rnn-to-transformer-lab
$ git checkout ch07
$ conda env create -f environment.yml
$ conda activate rnn-to-transformer-lab
$ python verify.py --only ch07
\end{minted}

Dòng cuối phải in ra \code{verify: ok}. Nếu không, đừng làm tiếp: mọi con số bạn
đo sau đó sẽ không so được với sách.

Một cảnh báo về cách chạy. Đừng gọi các script này
qua \code{conda run}: lệnh ấy gom stdout của tiến trình con rồi in lại bằng
encoding của hệ thống, nên output tiếng Việt làm nó ném
\code{UnicodeEncodeError} từ bên trong chính conda, và traceback không nhắc tới
file nào của repo. Gọi thẳng interpreter của env, hoặc activate env rồi chạy
\code{python}.

\begin{bt}
  \item Bỏ phép chia đi. Mở hàm \code{scaled\_dot\_product} trong
    \code{transformer.py} và xóa phép chia cho \code{math.sqrt(d\_k)}, rồi
    chạy \code{python -m pytest -q tests/test\_ch07.py} và nói test nào hỏng.
    Sau đó chạy lại \code{ch07\_corpus.py} và ghi lại cột \code{exact} ở cả bốn
    mốc epoch. \(\dmodel = 24\) và 4 head cho \(\dk = 6\); dựa vào bảng ở
    mục~\ref{sec:ch07-chia-can}, dự đoán trước khi chạy xem chỗ mất sẽ lớn hay
    nhỏ, rồi đối chiếu. Kết quả có mâu thuẫn với lập luận của mục đó không?

  \item Đổi hằng số \code{PE\_BASE} trong \code{transformer.py} từ 10000 thành
    100, rồi chạy lại script \code{ch07\_position.py}. Năm bảng của nó đổi thế
    nào? Với mỗi bảng, nói trước khi nhìn output xem nó \emph{phải} đổi hay
    \emph{phải} đứng yên, rồi kiểm. Sau đó chạy \code{ch07\_corpus.py} ở riêng
    mốc 14 epoch và nói con số ấy có đổi không.

  \item Cài lỗi thứ ba mà mục~\ref{sec:ch07-che} nêu: nhân mask vào sau softmax
    rồi \emph{quên} chuẩn hóa lại. Sửa \code{scaled\_dot\_product} cho đúng
    hình dạng ấy, chạy lại \code{ch07\_mask.py}, và trả lời hai câu. Phép can
    thiệp ở bảng đầu còn phát hiện được lỗi không? Mất mát huấn luyện đi theo
    chiều nào so với hai dòng đã có? Nói vì sao đây là lỗi khó thấy hơn cả lỗi
    bỏ mask.

  \item Lịch warmup của bài, làm cho tử tế. Mục~\ref{sec:ch07-corpus} thử nó ở
    \code{warmup} 100 và 400 với 658 bước huấn luyện và thua recipe dùng chung.
    Bây giờ cho nó đủ chỗ: chạy 56 epoch, tức khoảng 2632 bước, với công thức
    \(\dmodel^{-0.5}\min(s^{-0.5}, s \cdot w^{-1.5})\) ở vài giá trị \(w\). Có
    \(w\) nào đạt 0.9967 sớm hơn 42 epoch không? Nếu có, nói phát biểu nào của
    mục~\ref{sec:ch07-corpus} phải sửa lại.

  \item Quét số head. \code{ch07\_corpus.py} dựng đúng một Transformer ở hằng
    số \code{D\_MODEL} rồi chấm nó ở bốn mốc epoch, nên bài này phải bọc phần
    dựng và huấn luyện ấy trong một vòng lặp của riêng bạn: chạy \(h\) bằng 1,
    2, 4, 8 ở \(\dmodel = 24\), mỗi lần 42 epoch, và in số tham số cùng thời
    gian chạy cho từng \(h\). Kết quả có hình dạng giống hàng (A) của bảng 3
    không, và nếu khác thì corpus này thiếu cái gì mà WMT'14 có? Bốn lần huấn
    luyện vượt xa ngân sách 180 giây của script; chạy ngoài \code{verify.py} và
    nói rõ tổng thời gian bạn đã tiêu.
\end{bt}

\tierextend

\begin{bt}
  \item Mục~\ref{sec:ch07-chia-can} nói thẳng rằng nó chỉ đo lúc khởi tạo, và
    để ngỏ câu hỏi thật. Đo thống kê của \(\vect{q} \cdot \vect{k}\) trên một mô
    hình \emph{đã huấn luyện xong}: lấy điểm số thật ở từng tầng và từng head,
    rồi hỏi phương sai của chúng có còn tỷ lệ với \(\dk\) không, và tỷ số so với
    \(\dk\) là bao nhiêu. Nếu mô hình đã học ra một thang khác hẳn thang lúc
    khởi tạo, thì phép chia đang làm gì ở cuối quá trình huấn luyện? Kết luận
    của bạn phải kèm số ở ít nhất hai giá trị \(\dk\).

  \item Tách \enquote{thiếu bước} khỏi \enquote{thiếu dữ liệu}.
    Mục~\ref{sec:ch07-corpus} đo được rằng ba lần số epoch đưa Transformer tới
    chỗ mô hình attention tới ở 14 epoch, và đọc đó là chuyện thiên kiến quy
    nạp. Nhưng số epoch và số câu nhìn thấy dính vào nhau: 56 epoch trên 6000
    cặp là 336000 lượt nhìn. Thiết kế và chạy một thí nghiệm tách hai thứ, ví
    dụ giữ nguyên tổng số bước cập nhật mà đổi số cặp huấn luyện. Nếu 24000 cặp
    trong 14 epoch cho kết quả khác 6000 cặp trong 56 epoch, thì cái thiếu là
    dữ liệu; nếu giống, cái thiếu là bước.

  \item Đây là chỗ còn mở thật sự. \eqref{eq:ch07-tich-vo-huong-pe} cho một tích
    vô hướng chỉ phụ thuộc khoảng cách, và mục~\ref{sec:ch07-vi-tri} đo được nó
    không đơn điệu, nên khoảng cách không đọc ngược ra được. Câu hỏi: có bộ mã
    hóa vị trí nào mà \tn{tích vô hướng}{dot product} vừa chỉ phụ thuộc khoảng
    cách, vừa đơn điệu
    giảm trên toàn dải độ dài mà mô hình gặp không? Dựng một bộ như thế, chứng
    minh hai tính chất, rồi cắm vào repo và đo. Biết trước hai chuyện: bài đã đo
    ở hàng (E) rằng embedding học được cũng tốt ngang sinusoid, nên thanh chắn
    để vượt qua không phải sinusoid mà là \enquote{bất cứ thứ gì}; và
    chương~\ref{ch:vi-tri-va-ngu-canh} đọc dòng công trình đã đi theo hướng
    này, ở đó câu trả lời hay nhất tới nay không phải là cộng một bảng vào
    embedding.

  \item Bảng 1 của bài có ba cột và không có cột nào cho thứ
    mục~\ref{sec:ch07-corpus} đo được. Đề xuất một cột thứ tư đo \enquote{phải
    học bao nhiêu mới biết hai vị trí cạnh nhau}, định nghĩa nó đủ chặt để tính
    được cho cả ba hàng của bảng, rồi tính. Nếu bạn không định nghĩa được nó
    thành một đại lượng phụ thuộc riêng vào kiến trúc, hãy nói rõ nó phụ thuộc
    thêm vào cái gì, vì câu trả lời ấy chính là lý do bảng 1 không có cột đó.
\end{bt}
